\chapter{CNAEXT: Marker and Engine-Layer Catalog}
\label{ch:cnaext-catalog}
\label{ch:noxna-catalog}

Chapter~\ref{ch:design-philosophy}'s design-philosophy discussion introduced \texttt{NOXNA} as CNA's marker
convention --- and its matching compile-time purity check --- for any declaration inside the
XNA namespace tree that is not part of the real XNA~4.0 API. This appendix is not, and could
not responsibly be, an exhaustive line-by-line listing of every \texttt{NOXNA}-tagged member
in a codebase of several hundred header files; that list changes with every commit, and the
project's own \texttt{NOXNA} build gate (Chapter~\ref{ch:design-philosophy}) already checks it mechanically, which a
static list in a book cannot. What this appendix collects instead is every \texttt{NOXNA} (or
the closely related \texttt{EXT}-suffixed) extension named specifically, by class and member,
across this entire book --- organized by area, as a reference for the extensions this book
itself considers significant enough to have explained in the main text, with a pointer back
to where each one is covered in full.

\section{Core framework}

\begin{description}
\item[\cnaclass{Game}] \texttt{RunApplication},
  \texttt{getTargetFPS\allowbreak{}Property()}/%
  \texttt{getTargetMsFrameTime\allowbreak{}Property()}, static
  \texttt{fpsToMilliseconds\allowbreak{}PerFrame()}, and the framework-reserved F9/F10
  context-loss triggers in the ordinary event pump --- Ch.\ref{ch:game-loop}, Ch.\ref{ch:backend-architecture}.
\item[\cnaclass{GameWindow}] \texttt{GetNativeSdlWindowEXT()}, \texttt{IsBorderlessEXTProperty},
  \texttt{MinimizeEXT()} / \texttt{RestoreEXT()} --- Ch.\ref{ch:game-loop}.
\item[\cnaclass{GraphicsDeviceManager}] the entire \texttt{PresentationMode} enum
  (\texttt{Letterbox} / \texttt{Overscan} / \texttt{Stretch} / \texttt{NativeBackBuffer}/%
  \texttt{FixedHeightDynamicWidth}) and its get/set accessors --- Ch.\ref{ch:game-loop}.
\item[\cnaclass{Matrix}] \texttt{ToColumnMajor(float[16])}, the one graphics-motivated bridge
  method on an otherwise pure math type --- Ch.\ref{ch:math-core-types}.
\item[\cnaclass{Color}] two byte-based constructors and a commutative \texttt{float * Color}
  operator, both explicitly marked as conveniences absent from real XNA --- Ch.\ref{ch:math-core-types}.
\end{description}

\section{Graphics}

\begin{description}
\item[\cnaclass{Effect}] \texttt{Apply()} itself is \texttt{NOXNA} --- real FNA has no public
  \texttt{Effect.Apply()}, only \texttt{EffectPass.Apply()} --- Ch.\ref{ch:stock-effects}.
\item[\cnaclass{SkinnedEffect}] \texttt{VertexColorEnabled}, added for glTF-imported meshes
  carrying a \texttt{COLOR\_0} attribute; no such property exists on real XNA's own
  \cnaclass{SkinnedEffect} --- Ch.\ref{ch:stock-effects}.
\item[\cnaclass{GraphicsDevice}] \texttt{SetContextRecoveryEnabled()} and
  \texttt{SetStringMarkerEXT()}: the former combines a shared texture-shadow policy with an
  optional backend hook; the latter is native only on Vulkan, throws on D3D9, and is silent
  elsewhere --- Ch.\ref{ch:graphicsdevice}, Ch.\ref{ch:backend-architecture}.
\item[\cnaclass{IGraphicsBackend}]
  \texttt{SetStringMarkerEXT()} and \texttt{UpdatePresentationFormatEXT()}, plus the recovery
  setter and the two debug loss/restore hooks --- Ch.\ref{ch:backend-architecture}.
\item[Custom shaders] The entire \texttt{ShaderEffect} runtime-HLSL-compilation path on the
  three Direct3D backends --- Ch.\ref{ch:shader-fx-gap}, Ch.\ref{ch:direct3d-backends}.
\item[\cnaclass{PbrEffect}] An entire NOXNA glTF~2.0-style metallic-roughness PBR material
  effect with no real-XNA counterpart at all (\texttt{NormalMap}, \texttt{%
  MetallicRoughnessMap}, \texttt{EmissiveMap}, \texttt{OcclusionMap},
  \texttt{MetallicFactor} / \texttt{RoughnessFactor} / \texttt{EmissiveFactor}) --- Ch.\ref{ch:webgpu-backend}.
\item[\cnaclass{SkinnedPbrEffect}] PBR's separately-declared GPU-skinned sibling, rather
  than a mode on \cnaclass{PbrEffect}: it takes at most 72 bone transforms and requires the
  stride-68 tangent-equipped skinned-vertex layout (position, normal, tangent, UV, four weights
  and four byte indices). As with \cnaclass{SkinnedEffect}, game
  code supplies \texttt{AnimationPlayer::GetSkinTransforms()} to \texttt{SetBoneTransforms()}
  every frame; the loader does not do that binding. This is entirely \texttt{NOXNA}: real XNA
  predates both PBR and this PBR-plus-skinning combination. CNA keeps the class split in the
  same style as XNA's own distinct \cnaclass{BasicEffect} and \cnaclass{SkinnedEffect} types.
  Every backend with a real \cnaclass{PbrEffect} shader implements this sibling except WebGPU,
  whose current PBR path remains unskinned --- Ch.\ref{ch:models-meshes}, Ch.\ref{ch:webgpu-backend}.
\item[\cnaclass{AnimationPlayer}] The entire class, mirroring Microsoft's own unshipped XNA
  ``Skinned Model Sample'' reference code rather than anything in the real framework assembly
  --- \texttt{StartClip} / \texttt{Update} / \texttt{GetBoneTransforms}/\%
  \texttt{GetWorldTransforms} / \texttt{GetSkinTransforms}. Its companion
  \cnaclass{SkinningData}, plus \texttt{Keyframe} and \texttt{AnimationClip} aliases, are
  likewise \texttt{NOXNA} sample-derived types: the content reader attaches a
  \cnaclass{SkinningData} instance to \cnaclass{Model}'s \texttt{Tag}, preserving the original
  sample convention rather than inventing a new property on the faithful \cnaclass{Model}
  surface --- Ch.\ref{ch:models-meshes}.
\item[\cnaclass{MorphTargetEXT}] The entire glTF-style morph-target blending system ---
  \texttt{MorphTargetDataEXT}, \texttt{BlendMorphTargetsEXT}, \texttt{SetMorphWeightsEXT},
  \texttt{EvaluateMorphWeightsEXT} --- independent of \cnaclass{AnimationPlayer}'s own
  bone-track timeline, attached via a loaded \cnaclass{ModelMeshPart}'s own \texttt{Tag} ---
  Ch.\ref{ch:models-meshes}.
\item[\cnaclass{Texture}] Two static format-utility methods with no real-XNA counterpart,
  \texttt{GetBlockSizeSquaredEXT(SurfaceFormat)} and \texttt{GetFormatSizeEXT(SurfaceFormat)}
  --- the base class every texture type shares, used internally to compute OpenGL pixel-store
  alignment and to validate \texttt{SetData} / \texttt{GetData} byte counts --- Ch.\ref{ch:textures-rendertargets}.
\item[\cnaclass{TextureCube}] \texttt{DDSFromStreamEXT(GraphicsDevice\&, Stream\&)}, a static
  factory decoding a DDS-encoded cube map directly from a stream --- real XNA's own content
  pipeline never offered a runtime DDS loader for cube maps at all --- Ch.\ref{ch:textures-rendertargets}.
\item[\cnaclass{Texture3D}] \texttt{SetDataPointerEXT()}, a raw-pointer overload of the
  volume-texture upload path, alongside the already-documented finding (Ch.\ref{ch:content-and-assets}) that this same
  class's compressed-volume \texttt{byteCount}-vs-dimension validation was once dormant --- Ch.\ref{ch:textures-rendertargets}.
\item[\cnaclass{SurfaceFormat}] Seven enum values with no real-XNA equivalent ---
  \texttt{ColorBgraEXT}, \texttt{ColorSrgbEXT}, \texttt{Dxt5SrgbEXT}, \texttt{Bc7EXT},
  \texttt{Bc7SrgbEXT}, \texttt{ByteEXT}, \texttt{UShortEXT} --- each one recognized by
  \cnaclass{Texture}'s own \texttt{GetBlockSizeSquaredEXT} / \texttt{GetFormatSizeEXT} above ---
  Ch.\ref{ch:textures-rendertargets}.
\item[\cnaclass{GraphicsDevice}] \texttt{SetGraphicsProfileEXT(GraphicsProfile)}, an internal-use
  method letting \cnaclass{GraphicsDeviceManager} retroactively upgrade the eagerly-constructed
  default device (hardcoded \texttt{GraphicsProfile::Reach} at \cnaclass{Game} construction,
  before a game's own \texttt{GraphicsProfile} request can even be read) --- not exposed as a
  general runtime profile switch, since real XNA has none either --- Ch.\ref{ch:graphicsdevice}.
\item[\cnaclass{PresentationParameters}] \texttt{getHeadlessEXTProperty()} /
  \texttt{setHeadlessEXTProperty()} --- an opt-in, no-window device-creation mode for backends
  that can genuinely operate without a swap chain (today: D3D12 only; D3D11 and EasyGL throw,
  since their device construction is inseparable from a real window). This is the same
  no-window operating mode the \texttt{Headless} and \texttt{Software} backends already use
  unconditionally, offered here as an opt-in for a backend that normally does want one --- the
  mechanism this book's own screenshot infrastructure (see \texttt{tools/cna-screenshot-infra/}
  in this repository) and Ch.\ref{ch:verification-methodology}'s verification methodology both depend on --- Ch.\ref{ch:graphicsdevice}.
\end{description}

\section{Audio (Chapter~\ref{ch:audio-system})}

\cnaclass{DynamicSoundEffectInstance}'s float32 submission path ---
\texttt{SubmitFloatBufferEXT(const std::vector<float>\&)} and its offset/count overload --- has
no real-XNA counterpart at all; the faithful \texttt{SubmitBuffer} overloads only ever accepted
16-bit PCM byte buffers. The two paths are mutually exclusive once an instance is playing or
paused: submitting a float buffer switches the instance to float mode, and a subsequent
int16-mode \texttt{SubmitBuffer} call throws \texttt{System::InvalidOperationException} rather
than silently feeding raw int16 bytes into a live float-format stream, and symmetrically in the
other direction --- a \texttt{NOXNA} safety guard protecting a code path real XNA never had to
guard against.

\section{Media (Chapter~\ref{ch:media-system})}

\cnaclass{Video} and \cnaclass{VideoPlayer} share three \texttt{EXT}-suffixed members ---
\texttt{SetAudioTrackEXT(intcs)}, \texttt{SetVideoTrackEXT(intcs)}, and
\texttt{Video::FromUriEXT(std::string, GraphicsDevice*)} --- worth noting for a reason that
cuts against this appendix's usual grain: their own header comments describe each one as ``an
FNA extension beyond the original XNA~4.0 API surface,'' not a CNA invention. CNA's
\texttt{NOXNA} marker and \texttt{EXT} suffix convention is applied faithfully to code ported
from FNA's own multi-track-selection and URI-based video loading extensions, which already used
this identical naming pattern upstream --- not every entry in this appendix originates in this
project.

\section{Input (Chapter~\ref{ch:input-system})}

The largest concentration of \texttt{NOXNA} / \texttt{EXT} surface in the whole project. Whole
\texttt{NOXNA} subsystems: \cnaclass{CNA::Input::Clipboard}, \cnaclass{Joysticks},
\cnaclass{Sensors}, \cnaclass{Power}, \cnaclass{Haptics}, \cnaclass{InputDevices}. Per-class
\texttt{EXT} additions: \cnaclass{Keyboard}'s \texttt{GetKeyFromScancodeEXT}/%
\texttt{GetModStateEXT} / \texttt{GetScancodeNameEXT}; \cnaclass{Mouse}'s
\texttt{SetCursor} / \texttt{ClickedEXT} / \texttt{SetCaptureEXT} / \texttt{GetGlobalPositionEXT};
the entire \cnaclass{MouseCursor} class; \cnaclass{GamePad}'s \texttt{GetGUIDEXT}/%
\texttt{SetLightBarEXT} / \texttt{SetTriggerVibrationEXT} / \texttt{GetGyroEXT}/%
\texttt{GetAccelerometerEXT} / \texttt{GetPlayerIndexEXT} / \texttt{GetPowerInfoEXT}; and
\texttt{TextInputEXT}'s IME-candidate and input-type-hint additions.

\section{Devices and Sensors (Chapter~\ref{ch:devices-sensors})}

The entire \cnans{CNA::Devices} namespace is \texttt{NOXNA} in spirit --- \cnaclass{Camera},
\cnaclass{SystemTray}, \cnaclass{DisplayInfo}, \cnaclass{MessageBox}, \cnaclass{FileDialog},
\cnaclass{UrlLauncher}, \cnaclass{Locale}, \cnaclass{Clipboard}, \cnaclass{SystemInfo},
\cnaclass{PowerInfo} --- since none has any real XNA/WP7 precedent at all. Within the
faithful \cnans{Microsoft::Devices::Sensors} namespace: \cnaclass{SensorBase}'s
\texttt{TimeBetweenUpdatesChanged} event (a real, corrected mistagging finding --- it was
briefly and incorrectly marked as real XNA API before being fixed); \cnaclass{Compass}'s
\texttt{getStateProperty()} (added purely for symmetry with \cnaclass{Accelerometer}'s real
\texttt{State} property, which real XNA's \cnaclass{Compass} itself lacks);
\cnaclass{VibrateController}'s \texttt{Start(TimeSpan, float intensity)},
\texttt{getIsSupportedProperty()}, \texttt{getDeviceNameProperty()}, and
\texttt{StartLeftRight()}.

\section{Storage (Chapter~\ref{ch:storage})}

\cnaclass{StorageDevice}'s \texttt{SetAppNameEXT()} / \texttt{GetStorageRootEXT()} --- the two
methods the entire namespace's real, local-filesystem-backed behavior ultimately roots
itself in.

\section{Content and assets (Chapter~\ref{ch:content-and-assets})}

The \texttt{.cnj} content format itself is a \texttt{NOXNA} system in its entirety, alongside
the Avatar system's separately-kept \cnaclass{SkinnedModelEXT} content pipeline
(\texttt{.skinnedmodel.json} / \texttt{.skeleton.bin} / \texttt{.clip.bin}) and its own
\texttt{NOXNA} vertex/track/clip types (\cnaclass{VertexPositionNormalTextureSkinned},
\cnaclass{KeyframeEXT}, \cnaclass{BoneTrackEXT}, \cnaclass{AnimationClipEXT}) --- see
Chapter~\ref{ch:avatar}. \cnaclass{ContentManager} itself adds three further \texttt{NOXNA} methods purely
for build-time content auditing --- \texttt{RefreshContentManifest}, \texttt{%
GetContentManifest}, \texttt{GetXnbReaderUsageSummary} --- letting a game (or this book's own
research) enumerate every asset under a content root and which real \texttt{.xnb} reader name
actually loads it.

\section{Avatar (Chapter~\ref{ch:avatar})}

The entire \cnaclass{SkinnedModelEXT} real-rendering extension is \texttt{NOXNA} in its
entirety, deliberately opt-in and structurally separate from the faithful base Avatar API it
sits alongside. It includes \texttt{AvatarRenderer}'s \texttt{EnableRealRenderingEXT()} and
\texttt{DrawRealEXT()}, \texttt{SetAppearanceEXT()}, \cnaclass{AvatarAnimation}'s
\texttt{SetRealClipNameEXT} / \texttt{GetRealClipNameEXT} pair, and the CNA-invented
\cnaclass{AvatarAppearanceEXT} tint struct.

\section{GamerServices (Chapter~\ref{ch:gamerservices})}

\cnaclass{Guide}'s two real overlay-UI render hooks,
\texttt{RenderPendingMessageBoxEXT()} / \texttt{RenderPending\allowbreak{}Keyboard\allowbreak{}InputEXT()}, and their
matching headless test-simulation hooks. Beyond \cnaclass{Guide} itself: \cnaclass{LeaderboardEntry}'s
\texttt{getRankingEXTProperty()} (FNA's own extension --- its reference source already names
the field \texttt{RankingEXT} for the same reason as Ch.\ref{ch:media-system}'s video-track members above) and
\texttt{SetOnRatingChangedHookEXT()}, a CNA-added hook letting a rating change on a local-store
entry notify a caller (a no-op unless the hook is installed); \cnaclass{LeaderboardReader}'s
\texttt{ResliceEntriesEXT()}, the internal re-pagination step every read operation funnels
through; and \cnaclass{GamerPresence}'s \texttt{SetPresenceModeStringEXT()}.

\section{Networking (Chapter~\ref{ch:networking})}

\cnaclass{NetworkGamer}'s \texttt{SetHasLeftSession(bool)} is a real, and unusually
well-documented, case of this project's own ``restore a stubbed original member'' pattern (the
same shape as \cnaclass{Compass}'s \texttt{getStateProperty()} above, from the opposite
direction): FNA's own \texttt{HasLeftSession} setter is \texttt{private}, and FNA's own
\texttt{NetworkSession} never actually calls it after construction --- so on real, faithfully-ported
XNA/FNA semantics alone, \texttt{HasLeftSession} would be permanently \texttt{false} in
practice, an unimplemented stub rather than a functioning property. \cnaclass{NetworkGamer}
and \cnaclass{NetworkSession} are sibling classes, not base/derived, so CNA's
\cnaclass{NetworkSession::RemoveGamer()} could not have reached a real \texttt{private} setter
even if one were faithfully preserved; the header documents this reasoning directly rather than
silently exposing the setter. Marking it \texttt{NOXNA} and making it public is what lets this
port's own \texttt{RemoveGamer()} make \texttt{HasLeftSession} actually functional, unlike the
member it is restoring.

\section{How to read this appendix}

Treat this list the same way this book has asked you to treat every other list of specifics
it has offered: as a map into the source, not a replacement for it. The project's own
\texttt{NOXNA} compile-time gate (build with the \texttt{NOXNA} CMake option enabled) is the
authoritative, mechanically-enforced check for whether a given declaration inside the XNA
namespace tree is correctly tagged --- this appendix exists to give a reader who has just
finished this book a single place to recall which specific extensions it found significant
enough to explain, not to serve as a substitute for that build-time check against the
codebase's current, ever-changing state.
